\scanpage{218}
\deleted{\unsure{est}} \deleted{La relation signifie donc que pour} Ainsi

\[
\begin{aligned}
&\Pi\wedge_Z(S!)\xrightarrow{\sim} S^\vee\ \text{commute à un }g\in\unsure{\mathcal G|G}\ \text{donné}\\
&\text{si et seulement si pour tout }\varphi\in\Pi\ \text{-- ou}\\
&\text{simplement pour \uline{un} }\varphi_0\in\Pi\ \text{-- on a}\\
&\varphi(g_S(s))=g_{S^\vee}\varphi(s),\ \text{i.e. }\varphi\text{ commute à }g.
\end{aligned} \tag*{(78)}
\]

Donc, pour un sous-groupe donné $G\subset\mathcal G\subset\operatorname{Aut}(\Gamma)$ (on s'intéressera au cas où $Z\subset G$), l'isom. (76) est un isom de $G$-ens. ssi pour \uline{tout} $\varphi\in\Pi$ -- ou pour un $\varphi_0\in\Pi$, c'est kif-kif -- $\varphi$ commute à l'action de $\mathcal G$. \deleted{Le plus} Il y a un plus grand tel groupe $G$, qui \deleted{sont formé des $g\in\mathcal G$ qui} \inserted{est donné par}

\[
\mathcal G_0=\left\{g\in\operatorname{Aut}(S,S^\vee,R)\ \middle|\
\begin{aligned}
&\text{a) }g\text{ centralise }Z\ \text{i.e. }g\in\mathcal G\\
&\text{b) }\deleted{\unsure{\forall}}\exists\varphi\in\Pi,\ \varphi:S\xrightarrow{\sim}S^\vee\ \text{commute}\\
&\hspace{12mm}\text{à }g\ \text{i.e. }\varphi(g_Ss)=g_{S^\vee}(\varphi(s))
\end{aligned}\right\} \tag*{(79)}
\]

\margin{\uline{NB} Se garder de confondre cette condition de commutation au sens strict, avec la condition $g^\vee\varphi=\varphi g^{-1}$ (au lieu de $\varphi g$) !}

% slanted
\margin{Exemple un peu idiot, pas très instructif. Exemple important par contre : celui de $S,S^\vee$ déduits de $M^*,M^{\vee *}$ (p. 104' ff) \deleted{\unsure{alors}} Alors la \unsure{partie|trace} de $\mathcal G_0$ sur le groupe $\operatorname{Gl}_{\Lambda}^{\unsure{\prime|{}}}(M)$ $\mathcal G_0$ est formé des \unsure{il faut|u qui} $u\in\operatorname{Gl}_\Lambda(M)$ satisfaisant $\check u\varphi_0=\pm\varphi_0u$ i.e. $\underbrace{\check u\varphi_0u^{-1}}_{\varphi_0\det(u)}=\pm\varphi_0$ i.e. $\det u\in\pm1$. On a alors $(79\ \mathrm{bis})\quad\mathcal G_0=\operatorname{Gl}'_\Lambda(M)^{\pm1}$ du moins la \unsure{trace|prise|partie} du $\mathcal G_0$ \unsure{l'y|(79)} sur $\operatorname{Gl}'_\Lambda(M)$ \ldots}

\uline{Ex} Prenons $Z=1$ donc $\mathcal G=\operatorname{Aut}(\Gamma)$. Alors $\Pi$ se réduit à un élément, qui est n'importe quelle jauge $\varphi_0$, de $\Gamma=(S,S^\vee,R)$. Bien entendu, il n'est pas obligé que $\varphi_0$ commute à l'action de $\mathcal G$ -- p. ex. dans le cas de la p. 110 $S=S^\vee$, $R={}$ diagonale, $\varphi_0$ est donc une involution de $S$, tandis que $\mathcal G=\unsure{\mathfrak S|\mathscr S}_S$ ; ici $\mathcal G_0$ s'identifie au commutant de l'involution $\varphi_0$ dans $\mathfrak S_{E}$. \deleted{\unsure{Dans ce cas|Donc on a}}

Il faut encore voir comment se transforme par dualité une structure de multigraphe \inserted{(\deleted{\unsure{strict}} séparé)} à homothéties, \unsure{interprété|réinterprété} dans l'optique présente

% archive sheet number: 217
